\section{Discussion and Limitations}

The audit supports a qualified use of post-hoc explanations: they can provide
reproducible accounts and measurable evidence about frozen models, provided the
account's target, identifiable support and evaluation rule are specified.
Reproducibility does not imply that an intervention will outperform a matched
random subset. Conversely, an intervention response describes the model under
that replacement mechanism, not a feature's causal role in an attack.
The practical reporting unit should therefore retain the method, target,
reference, operator, panel and uncertainty together. Collapsing these into a
trust score would hide distinctions that the experiments expose.

\emph{Computational scope.} A cached two-input CICIDS/MLP microbenchmark on an
Apple M1 Pro CPU (10 cores, 32 GiB), excluding the first repetition and
preparation, measures one explanation of one input per observation, not a
joint two-input batch. Table~\ref{tab:binary-cost} reports the median, range
and measurement count.
\begin{table}[t]
\centering\small
\caption{Single-input CPU timings (ms); first repetition excluded.
Two fixed inputs supply the $n$ individual measurements per method.}
\label{tab:binary-cost}
\begin{tabular}{lrrr}
\toprule
Method & Median & Min--max & $n$ \\
\midrule
G$\times$I & 4.4 & 3.7--4.6 & 8 \\
IG & 238.5 & 233.5--256.9 & 8 \\
Occlusion & 166.3 & 163.3--191.4 & 8 \\
LIME & 171.8 & 160.2--174.6 & 6 \\
SHAP & 988.3 & 899.5--1,101.1 & 6 \\
\bottomrule
\end{tabular}

\end{table}
These are small-sample descriptive timings: gradient/Occlusion batches were
one, and LIME/SHAP prediction batches 4,096. SHAP used one physical background
row, unlike the main audit's 100 copies. One explanation is not LIME's
50-run recomputation experiment or the full intervention audit. Historical
batch-stage costs include preparation and are not pooled with these timings.
Full-audit critical-path time and operational per-alert latency remain unknown.

\emph{Limits.} One fixed split and five fitted seeds do not quantify new-site,
temporal or session-level generalisation. Exact/rounded overlap safeguards
do not establish complete independence; source identity is also not a
session cluster. Class caps alter prevalence, and the correct-by-both panel
conditions RQ1/RQ2 on successful decisions. The processed ToN-IoT extract lacks
timestamps and endpoint addresses needed for stronger split policies.
Feature-space replacements need not describe feasible traffic. The CNN's
feature order and selected configurations limit architectural interpretation.

The RQ structure and operator/outcome extensions reorganise existing
experiments; sampling budgets were increased after inspecting earlier
results. They are not retroactively preregistered hypotheses. SHAP's two-run
support and unknown historical regularisation/environment further limit
reproducibility claims. The explicitly configured SHAP 0.51.0 microbenchmark
does not recover the historical runs' settings. Marginal intervals and
exploratory example selection do not establish a universal method winner.
The study measures neither analyst understanding/trust nor detection gains
from using explanations, and proposes no operational safety threshold.

\section{Conclusion}

For the evaluated binary tasks and fitted pipelines, post-hoc attributions
provide condition-specific evidence about model behaviour, but do not
independently validate the correctness or safety of IDS decisions.
Identifiable accounts can be highly reproducible while
functional conclusions change with the replacement operator. Outcome diagnostics
expose further dependence on the model, score and available support. An
auditable explanation must therefore state these conditions; explaining a
decision is not equivalent to validating an IDS.
